% Part: history
% Chapter: biographies
% Section: julia-robinson
\documentclass[../../../include/open-logic-section]{subfiles}

\begin{document}

\olfileid{his}{bio}{rob}

\olsection{جولیا رابینسون}

\olphoto{robinson-julia}{Julia Robinson}

جولیا بومن رابینسون ریاضی‌دانی آمریکایی بود. او بیشتر به سبب کارش روی
مسئله‌های تصمیم، و مشهورتر از همه سهمش در حل مسئلهٔ دهم هیلبرت،
شناخته می‌شود. رابینسون در 8 دسامبر 1919 در سنت لوئیس میزوری به دنیا
آمد. او به یاد می‌آورد که از کودکی شیفتهٔ اعداد بود
\citep[4]{Reid1986}. در 9سالگی به مخملک مبتلا شد و چند بار نیز تب
روماتیسمی او عود کرد. ازاین‌رو ناچار بود زمان زیادی را در بستر بگذراند
و از تحصیل عقب افتاد. هرچند با کمک آموزگاران خصوصی توانست عقب‌ماندگی
را جبران کند، آثار جسمانی بیماری تا پایان عمر بر او باقی ماند.

رابینسون با وجود دشواری‌های کودکی، دبیرستان را با چند جایزه در ریاضیات
و علوم به پایان رساند. تحصیلات دانشگاهی را در کالج ایالتی سن‌دیگو آغاز
کرد و در سال آخر به دانشگاه کالیفرنیا در برکلی منتقل شد. آنجا از
رافائل رابینسونِ ریاضی‌دان تأثیر پذیرفت. آن دو دوستان خوبی شدند و در
1941 ازدواج کردند. جولیا به‌عنوان همسر یکی از اعضای هیئت علمی اجازه
نداشت در گروه ریاضی برکلی تدریس کند. گرچه همچنان آزادانه در کلاس‌های
ریاضی حاضر می‌شد، امیدوار بود دانشگاه را ترک کند و خانواده تشکیل دهد.
اما اندکی پس از ازدواج به ذات‌الریه مبتلا شد. به او گفتند تب روماتیسمی
دوران کودکی موجب تجمع چشمگیر بافت جوشگاهی روی قلبش شده است. پزشک به
سبب شدت این بافت پیش‌بینی کرد که بیش از چهل سال زنده نخواهد ماند و
توصیه کرد فرزندی نداشته باشد \citep[13]{Reid1986}.

رابینسون مدتی طولانی افسرده بود، اما سرانجام تصمیم گرفت مطالعهٔ ریاضیات
را ادامه دهد. به برکلی بازگشت و در سال 1948 زیر نظر آلفرد تارسکی دکتری
گرفت. تارسکی نشان داده بود نظریهٔ مرتبهٔ اول اعداد حقیقی تصمیم‌پذیر
است و از کار گودل نتیجه می‌شد نظریهٔ مرتبهٔ اول اعداد طبیعی تصمیم‌ناپذیر
است. اینکه نظریهٔ مرتبهٔ اول اعداد گویا تصمیم‌پذیر است یا نه مسئلهٔ
باز مهمی بود. رابینسون در رسالهٔ خود \citeyearpar{Robinson1949} اثبات
کرد که تصمیم‌پذیر نیست.

رابینسون که به مسئله‌های تصمیم علاقه‌مند بود، سپس کوشید راه‌حلی برای
مسئلهٔ دهم هیلبرت بیابد. این مسئله یکی از فهرست مشهور 23 مسئلهٔ ریاضی
بود که داوید هیلبرت در سال 1900 مطرح کرد. مسئلهٔ دهم می‌پرسد آیا
الگوریتمی وجود دارد که در زمانی متناهی تعیین کند یک معادلهٔ چندجمله‌ای
با ضریب‌های صحیح، مانند $3x^2 - 2y +3 = 0$، در اعداد صحیح جواب دارد یا
نه. چنین پرسش‌هایی \emph{مسئله‌های دیوفانتی} نامیده می‌شوند. رابینسون
پس از چند موفقیت اولیه با مارتین دیویس و هیلاری پاتنم، که آنان نیز روی
مسئله کار می‌کردند، همراه شد. آن‌ها نشان دادند مسئله‌های دیوفانتی
نمایی، که در آن‌ها مجهول‌ها می‌توانند در توان نیز ظاهر شوند،
تصمیم‌ناپذیرند و نشان دادند حدسی معین، که بعدها «J.R.» نام گرفت،
تصمیم‌ناپذیری مسئلهٔ دهم هیلبرت را نتیجه می‌دهد
\citep{DavisPutnamRobinson1961}. رابینسون در سراسر دههٔ 1960 به کار روی
این مسئله ادامه داد. در سال 1970 یوری ماتیاسویچ، ریاضی‌دان جوان روس،
سرانجام فرضیهٔ J.R. را اثبات کرد. نتیجهٔ ترکیبی اکنون قضیهٔ
ماتیاسویچ--رابینسون--دیویس--پاتنم، یا به اختصار قضیهٔ MRDP، نامیده
می‌شود. ماتیاسویچ و رابینسون دوست شدند و در چند مقاله همکاری کردند.
رابینسون یک بار در نامه‌ای به ماتیاسویچ نوشت: «واقعاً بسیار خرسندم که
با کار مشترک، با وجود هزاران کیلومتر فاصله، آشکارا بیش از هر یک از ما
به‌تنهایی پیشرفت می‌کنیم» \citep[45]{Matijasevich1992}.

رابینسون نخستین رئیس زن انجمن ریاضی آمریکا و نخستین زنی بود که به
عضویت آکادمی ملی علوم انتخاب شد. او پس از تشخیص سرطان خون، در
30 ژوئیهٔ 1985 و در 65سالگی درگذشت.

\begin{reading}
مقاله‌های ریاضی رابینسون در \textit{Collected Works} او در دسترس‌اند
\citep{Robinson1996}. این مجموعه بازچاپ یادنامهٔ زندگی‌نامه‌ای او در
آکادمی ملی علوم را نیز در بر دارد \citep{Feferman1994}. کنستانس رید،
خواهر بزرگ‌تر رابینسون، «زندگی‌نامهٔ خودنوشت جولیا» را بر پایهٔ
گفت‌وگوها منتشر کرد \citep{Reid1986} و نیز یادنامه‌ای کامل نوشت
\citep{Reid1996}. مستندی کوتاه دربارهٔ رابینسون و مسئلهٔ دهم هیلبرت را
جرج چیسری کارگردانی کرده است \citep{Csicsery2016}. برای یادداشتی کوتاه
دربارهٔ همکاری‌های یوری ماتیاسویچ با رابینسون و تأثیر او بر کار
ماتیاسویچ، \citep{Matijasevich1992} را ببینید.
\end{reading}

\end{document}
